\section{Modelado e identificación del sistema}

\subsection{Modelo físico}

El conjunto formado por el transistor calentador y el sensor se representó como una única masa a
temperatura uniforme. Bajo esa hipótesis, y considerando pérdidas por convección natural y por
radiación hacia el entorno, el balance de energía conduce a \cite{hedengren_modeling}

\begin{equation}
m c_p \frac{dT}{dt} = \alpha Q - UA\,(T - T_a) - \epsilon\sigma A\,(T^4 - T_\infty^4),
\label{eq:nolineal}
\end{equation}

donde $\alpha$ convierte el porcentaje de potencia comandado en vatios entregados al sistema, $UA$
es el coeficiente global de convección, y $\epsilon$ y $\sigma$ son la emisividad y la constante de
Stefan--Boltzmann. Las hipótesis adoptadas fueron masa concentrada y temperatura uniforme,
convección predominantemente natural, potencia del calentador proporcional al comando, y
$T_a \approx T_\infty$ por suponer el equipo y el recinto inicialmente en equilibrio térmico. Las
temperaturas del término radiativo se manejaron en kelvin.

La no linealidad de \eqref{eq:nolineal} proviene exclusivamente del término $T^4$. Linealizando
alrededor de un punto de operación $(\bar T, \bar Q)$ y definiendo las variables de desviación
$\Delta T = T - \bar T$ y $\Delta Q = Q - \bar Q$, se obtiene \cite{hedengren_linearize}

\begin{equation}
\frac{d\,\Delta T}{dt} = \gamma\,\Delta T + \beta\,\Delta Q,
\label{eq:lineal}
\end{equation}

\vspace{-3mm}
\begin{equation}
\gamma = -\frac{UA + 4\epsilon\sigma A \bar T^{3}}{m c_p},
\qquad
\beta = \frac{\alpha}{m c_p},
\label{eq:gammabeta}
\end{equation}

cuya función de transferencia es de primer orden,

\begin{equation}
G(s) = \frac{K}{\tau s + 1},
\qquad \tau = -\frac{1}{\gamma},
\qquad K = -\frac{\beta}{\gamma}.
\label{eq:g1orden}
\end{equation}

En \eqref{eq:lineal} el coeficiente $\gamma$, dado por \eqref{eq:gammabeta}, reúne las dos vías de
pérdida, y es el término radiativo el que lo hace depender del punto de operación. El punto empleado fue
$\bar Q = \SI{50}{\percent}$; los valores de $\tau$ y $K$ que resultan de \eqref{eq:g1orden} se
resumen en la Tabla~\ref{tab:fisico}.

\begin{table}[!t]
\caption{Modelo físico linealizado en $\bar Q = \SI{50}{\percent}$ y error de la aproximación
lineal frente al modelo no lineal.}
\label{tab:fisico}
\centering
\footnotesize
\begin{tabular}{@{}lr@{}}
\toprule
\textbf{Magnitud} & \textbf{Valor} \\
\midrule
Temperatura de equilibrio $\bar T$        & \SI{60.47}{\celsius} \\
Pérdida por convección en equilibrio      & \SI{0.2188}{\watt} (\SI{43.8}{\percent}) \\
Pérdida por radiación en equilibrio       & \SI{0.2811}{\watt} (\SI{56.2}{\percent}) \\
$\gamma$                                  & \SI{-7.548e-3}{\per\second} \\
$\beta$                                   & \SI{5.000e-3}{\celsius\per\percent\per\second} \\
Constante de tiempo $\tau$                & \SI{132.5}{\second} \\
Ganancia estática $K$                     & \SI{0.6625}{\celsius\per\percent} \\
\midrule
\multicolumn{2}{@{}l}{\textit{Error del modelo lineal frente al no lineal}} \\
$\Delta Q = \SI{5}{\percent}$: error final / RMSE   & \SI{0.029}{\celsius} / \SI{0.022}{\celsius} \\
$\Delta Q = \SI{40}{\percent}$: error final / RMSE  & \SI{1.744}{\celsius} / \SI{1.337}{\celsius} \\
\bottomrule
\end{tabular}
\end{table}

Los dos últimos renglones de la Tabla~\ref{tab:fisico} delimitan el alcance de la aproximación
lineal. Ante una perturbación de \SI{5}{\percent} el error no llega a \SI{0.03}{\celsius}, mientras
que al multiplicar la amplitud por ocho el error crece casi sesenta veces. Este crecimiento no es
proporcional porque el término omitido es de cuarto orden en la temperatura, y en el punto de
operación la radiación ya aporta el \SI{56.2}{\percent} de las pérdidas totales. Es decir, el
modelo lineal es útil para excursiones pequeñas alrededor del punto de operación, pero no para
describir el calentamiento completo desde temperatura ambiente, que es justamente el transitorio
que debe manejar el controlador. Esta limitación motivó obtener el modelo de diseño de forma
experimental en lugar de usar los parámetros físicos.

\subsection{Identificación experimental}

Se realizaron dos experimentos sobre la unidad TCLab n.\textsuperscript{o}~4, con temperatura
ambiente de \SI{24}{\celsius} y periodo de muestreo $T_s = \SI{1}{\second}$. El primero aplicó un
escalón de $Q_1$ de \SIrange{0}{50}{\percent} y se registró durante \SI{1860}{\second}
(\num{1861} muestras), partiendo de $T_1 = \SI{24.78}{\celsius}$. El segundo aplicó una secuencia
multinivel con escalones sucesivos de \SIrange{10}{50}{\percent} durante \SI{3520}{\second}
(\num{3521} muestras), partiendo de $T_1 = \SI{25.27}{\celsius}$. El ensayo escalón concentra la
información en una sola transición amplia, mientras que el multinivel recorre varios niveles de
potencia y permite observar si la ganancia del proceso depende de la región de operación.

Ambos conjuntos se ajustaron a un modelo de primer orden con tiempo muerto (FOPDT)
\cite{hedengren_fopdt},

\begin{equation}
G_p(s) = \frac{K\,e^{-\theta s}}{\tau s + 1},
\label{eq:fopdt}
\end{equation}

expresado en variables de desviación, con $[K] = \si{\celsius\per\percent}$ y
$[\tau] = [\theta] = \si{\second}$. Los parámetros de \eqref{eq:fopdt} se obtuvieron por dos vías: el procedimiento de
mínimos cuadrados sobre la respuesta simulada empleado en la práctica, y la función
\texttt{procest} de \emph{System Identification Toolbox} con estructura \texttt{P1D}. La
Tabla~\ref{tab:fopdt} recoge los cuatro modelos y la Fig.~\ref{fig:experimentos} muestra las
señales experimentales junto con la respuesta de los modelos identificados.

\begin{table}[!t]
\caption{Parámetros FOPDT identificados por ensayo y por método.}
\label{tab:fopdt}
\centering
\footnotesize
\begin{tabular}{@{}llS[table-format=1.3]S[table-format=3.1]S[table-format=1.2]S[table-format=2.1]@{}}
\toprule
\textbf{Ensayo} & \textbf{Método} & {$K$ (\si{\celsius\per\percent})} & {$\tau$ (\si{\second})}
& {$\theta$ (\si{\second})} & {Ajuste (\si{\percent})} \\
\midrule
Escalón    & Práctica              & 0.979 & 179.1 & 8.00 & 93.0 \\
Escalón    & \texttt{procest} P1D  & 0.979 & 178.7 & 8.36 & 93.0 \\
Multinivel & Práctica              & 1.036 & 189.5 & 2.00 & 94.5 \\
Multinivel & \texttt{procest} P1D  & 1.036 & 189.2 & 2.80 & 94.5 \\
\bottomrule
\end{tabular}
\end{table}

\begin{figure*}[!t]
\centering
\begin{tikzpicture}
\begin{groupplot}[group style={group size=2 by 2, horizontal sep=1.3cm, vertical sep=0.9cm},
                  estilo, ancho]
\nextgroupplot[ylabel={$Q_1$ (\si{\percent})}, xticklabels={}, ymin=-5, ymax=60,
               title={\scriptsize Ensayo escalón}, title style={yshift=-2pt}]
  \addplot[cNaranja, aligera=3] table[col sep=comma, x=Tiempo_s, y=Q1_pct] {\G/IdentificacionSE.csv};
\nextgroupplot[ylabel={$Q_1$ (\si{\percent})}, xticklabels={}, ymin=-5, ymax=60,
               title={\scriptsize Ensayo multinivel}, title style={yshift=-2pt}]
  \addplot[cNaranja, aligera=5] table[col sep=comma, x=Tiempo_s, y=Q1_pct] {\G/IdentificacionSM.csv};
\nextgroupplot[ylabel={$T_1$ (\si{\celsius})}, xlabel={Tiempo (\si{\second})}, legend pos=south east]
  \addplot[cAzul, aligera=3] table[col sep=comma, x=Tiempo_s, y=T1_experimental_C] {\G/IdentificacionSE.csv};
  \addplot[cVerde, dashed, aligera=3] table[col sep=comma, x=Tiempo_s, y=T1_modelo_C] {\G/IdentificacionSE.csv};
  \legend{Experimental, FOPDT}
\nextgroupplot[ylabel={$T_1$ (\si{\celsius})}, xlabel={Tiempo (\si{\second})}, legend pos=south east]
  \addplot[cAzul, aligera=5] table[col sep=comma, x=Tiempo_s, y=T1_experimental_C] {\G/IdentificacionSM.csv};
  \addplot[cVerde, dashed, aligera=5] table[col sep=comma, x=Tiempo_s, y=T1_modelo_C] {\G/IdentificacionSM.csv};
  \legend{Experimental, FOPDT}
\end{groupplot}
\end{tikzpicture}
\caption{Señales de los dos experimentos de identificación y respuesta de los modelos FOPDT
ajustados a cada uno. Izquierda: escalón de $Q_1$ a \SI{50}{\percent}. Derecha: secuencia
multinivel de \SIrange{10}{50}{\percent}.}
\label{fig:experimentos}
\end{figure*}

Los dos métodos coinciden en $K$ y $\tau$ hasta la tercera cifra significativa, lo que respalda la
identificación. La única diferencia apreciable está en $\theta$: el procedimiento de la práctica
cuantiza el tiempo muerto a un número entero de periodos de muestreo, mientras que
\texttt{procest} lo estima como variable continua. Las diferencias resultantes, \SI{0.36}{\second}
y \SI{0.80}{\second}, son menores que un periodo de muestreo y por tanto no distinguibles con los
datos disponibles.

La coincidencia entre ambos métodos depende de entregar los datos en variables de desviación: con
las señales absolutas la estimación converge a $K \approx \SI{1.47}{\celsius\per\percent}$ y
$\theta = 0$, valor que no es la ganancia del proceso sino el cociente
$T_1(\infty)/\Delta Q_1$, porque el estimador debe explicar también el nivel inicial de
temperatura.

\subsection{Validación cruzada y selección del modelo}

Cada modelo se simuló con la entrada del experimento que no se usó para identificarlo, manteniendo
fijos $K$, $\tau$ y $\theta$. La Tabla~\ref{tab:validacion} reporta las métricas calculadas desde
el primer cambio de $Q_1$, que es el intervalo que evalúa la respuesta dinámica, y la
Fig.~\ref{fig:validacion} muestra el cruce del modelo seleccionado junto con su residuo.

\begin{table}[!t]
\caption{Validación cruzada. Métricas calculadas desde el primer cambio de $Q_1$.}
\label{tab:validacion}
\centering
\footnotesize
\begin{tabular}{@{}lS[table-format=1.2]S[table-format=1.2]S[table-format=+1.2]S[table-format=1.3]S[table-format=2.1]@{}}
\toprule
\textbf{Modelo $\rightarrow$ datos} & {RMSE} & {MAE} & {Sesgo} & {$R^2$} & {Ajuste} \\
 & {(\si{\celsius})} & {(\si{\celsius})} & {(\si{\celsius})} & & {(\si{\percent})} \\
\midrule
Escalón $\rightarrow$ multinivel & 1.62 & 1.51 & +1.50 & 0.987 & 88.7 \\
Multinivel $\rightarrow$ escalón & 2.63 & 2.41 & -2.41 & 0.934 & 74.4 \\
\bottomrule
\end{tabular}
\end{table}

\begin{figure}[!t]
\centering
\begin{tikzpicture}
\begin{groupplot}[group style={group size=1 by 2, vertical sep=0.7cm}, estilo, height=3.6cm]
\nextgroupplot[ylabel={$T_1$ (\si{\celsius})}, xticklabels={}, legend pos=south east]
  \addplot[cAzul, aligera=5] table[col sep=comma, x=Tiempo_s, y=T1_experimental_C] {\G/validacionEscEnMul.csv};
  \addplot[cVerde, dashed, aligera=5] table[col sep=comma, x=Tiempo_s, y=T1_modelo_C] {\G/validacionEscEnMul.csv};
  \legend{Experimental, Modelo escalón}
\nextgroupplot[ylabel={Residuo (\si{\celsius})}, xlabel={Tiempo (\si{\second})}]
  \addplot[cRojo, aligera=5] table[col sep=comma, x=Tiempo_s, y=Residuo_C] {\G/validacionEscEnMul.csv};
  \addplot[cGris, dashed, domain=0:3520, samples=2] {0};
\end{groupplot}
\end{tikzpicture}
\caption{Validación del modelo identificado con el ensayo escalón sobre los datos multinivel, y
residuo de predicción $e(k) = T_{1,\mathrm{exp}}(k) - T_{1,\mathrm{mod}}(k)$.}
\label{fig:validacion}
\end{figure}

El modelo del ensayo escalón predice claramente mejor los datos independientes: su RMSE es un
\SI{38}{\percent} menor y su ajuste es \num{14.3} puntos porcentuales mayor. Ninguno de los dos
cruces está libre de sesgo. El modelo escalón subestima la temperatura en promedio
\SI{1.50}{\celsius}, y el residuo de la Fig.~\ref{fig:validacion} crece de forma ordenada con el
nivel de potencia: \num{+1.29}, \num{+0.88}, \num{+1.23}, \num{+1.85} y \SI{+2.20}{\celsius} para
$Q_1$ de \num{10} a \SI{50}{\percent}. El modelo multinivel presenta el sesgo opuesto,
\SI{-2.41}{\celsius}, es decir, sobreestima la temperatura del ensayo escalón. Ambos patrones
apuntan al mismo origen: la ganancia efectiva de la planta no es constante, de modo que un modelo
identificado en una región tiende a errar sistemáticamente en la otra. Esto es coherente con lo
observado en la Sección~II-A, donde el peso creciente del término radiativo hacía que una única
recta no describiera todo el rango.

Dado que ninguno de los dos modelos es válido en todo el rango, el criterio de selección fue la
región de operación prevista para el controlador. La referencia de trabajo es \SI{60}{\celsius},
que corresponde aproximadamente a $Q_1 = \SI{50}{\percent}$, justamente la condición del ensayo
escalón. Ese modelo además tuvo mejor capacidad predictiva sobre datos independientes. Su tiempo
muerto, $\theta = \SI{8.00}{\second}$, equivale a ocho periodos de muestreo, lo que permite
representar el retardo de forma exacta en la implementación discreta. El modelo nominal adoptado es
entonces

\begin{equation}
\Gpnom(s) = \frac{0{,}979\;e^{-8{,}00\,s}}{179{,}1\,s + 1}
\quad \left[\si{\celsius\per\percent}\right].
\label{eq:gpnom}
\end{equation}
